\documentclass[10pt]{extarticle}

\usepackage{fancyhdr, geometry, hyperref, xcolor, amsmath, amsfonts, subcaption, graphicx, tikz, nomencl, bm}
\usetikzlibrary{shapes.geometric}
\usetikzlibrary{arrows.meta,arrows}
\usepackage[export]{adjustbox}
\usepackage[utf8]{inputenc}
\usepackage{array}
\usepackage{pdflscape}

\geometry {
    a4paper,
    total={170mm,257mm},
    left=20mm,
    top=20mm,
}

\pagestyle{fancy}
\pagenumbering{arabic}
\fancyhead[L]{Bike - Simplified}
\fancyfoot[C]{}
\fancyfoot[R]{\thepage}

\hypersetup{
    colorlinks,
    linkcolor={blue!50!black},
    citecolor={blue!50!black},
    urlcolor={blue!80!black}
}

\numberwithin{equation}{subsubsection}
\makenomenclature
\newcommand{\sgn}{\operatorname{sgn}}

\newcommand{\dr}[2]{$\displaystyle #1$ & $\displaystyle #2$ \\[4pt]}
\newcommand{\drl}[1]{$\displaystyle #1$ & \\[4pt]}

\newenvironment{twocol}{%
    \begin{center}
    \begin{tabular}{@{} l @{\qquad\quad} l @{}}
}{%
    \end{tabular}
    \end{center}
}

\begin{document}
\section{Bike systems}

\subsection{List}
\begin{enumerate}
    \item Frame
    \item Wheel rear
    \item Wheel front
    \item Drivetrain
\end{enumerate}
To Add
\begin{itemize}
    \item FBD, dimensions diagrams
\end{itemize}

% =====================================================================
\newpage
\subsection{Bike}
% =====================================================================

% COMMENT: "total bike system" means frame + both wheels + drivetrain treated as one body.
% The internal forces (axle reactions, chain tension) do not appear here — they cancel.
The equation of motion is written for the total bike system. Internal forces between subsystems cancel and do not appear.

\textbf{Assumptions}
\begin{itemize}
    \item Rigid frame --- no compliance, no rider motion relative to frame
    \item Point mass at fixed height $h$ --- rider and frame CG collocated and invariant
    \item Motion confined to the $x$--$z$ plane --- straight-line travel, no yaw or roll
    \item Constant road gradient $\theta_1$ --- no curvature or undulation
    \item Two contact points only --- front and rear tyre, no additional supports
    \item Aerodynamic lift neglected --- drag and pitching moment only
    \item Wind in the $x$--$z$ plane --- no crosswind
    \item Rolling resistance as contact resultant shift (Option A) --- no separate resistance moment
    \item Both tyres share the same $c_{rr}$ and $r_w$
\end{itemize}

\subsubsection{Forces}

% COMMENT: weight is in the ground frame here; the road-frame projection appears in Translation.
Weight acts at the system centre of gravity $a$ in the ground frame:
\begin{equation}
    \vec{W}_1 = mg\,\hat{z}_g
\end{equation}

% COMMENT: tire deformation shifts the contact resultant forward of the geometric contact point
% by e = c_rr * r_w. C and D are therefore not the geometric contact points but the shifted ones.
% This is Option A: the shift is absorbed into the moment arms in Geometry; no separate rolling
% resistance moment appears in the EOM.
Under load, tire deformation shifts the contact resultant forward of the geometric contact by $e = c_{rr}r_w$. Points $C$ and $D$ denote these shifted contact locations. The ground reactions at $C$ and $D$ are decomposed in the road frame $(\hat{x}_b, \hat{z}_b)$:
\begin{gather}
    \vec{R}_{2,0C} = R_{2,0Cx}\,\hat{x}_b + R_{2,0Cz}\,\hat{z}_b \\
    \vec{R}_{3,0D} = R_{3,0Dx}\,\hat{x}_b + R_{3,0Dz}\,\hat{z}_b
\end{gather}

% COMMENT: aerodynamic lift is neglected. A_{1m} is the pitching moment (nose-down positive here).
% Both are expressed in the road frame; wind direction enters through v_rel in the expansion.
Aerodynamic drag and pitching moment act on the frame, neglecting lift:
\begin{gather}
    \vec{A}_{1f} = -A_{1fx}\,\hat{x}_b \\
    \vec{A}_{1m} = -A_{1my}\,\hat{y}_b
\end{gather}

\subsubsection{Geometry}

% COMMENT: l_2, l_3 are measured from the CG to the geometric wheel contact points (not the shifted ones).
% The effective arms l_2', l_3' absorb the rolling resistance offset e*sigma and are used in the moment equation.
% sigma = sgn(x_dot) keeps the offset directed forward regardless of travel direction;
% for a bike sigma = +1 always but it is kept for formal completeness.
% aA is the vector from CG to the rear axle, used as the aerodynamic pitching moment application point.
Horizontal distances $l_2$, $l_3$ are measured from the system CG $a$ to the geometric wheel contact points. The rolling resistance offset $e$ shifts the contact resultant forward; $\sigma$ ensures this shift is always in the direction of motion:
\begin{gather}
    L = l_2 + l_3 \qquad e = c_{rr}\,r_w \qquad \sigma = \sgn(\dot{x}_1) \\
    l_2' = l_2 - e\sigma \qquad l_3' = l_3 + e\sigma
\end{gather}
The position vectors from $a$ to each force application point, expressed in the road frame:
\begin{gather}
    \vec{aC} = (-l_2 + e\sigma)\,\hat{x}_b + h\,\hat{z}_b \\
    \vec{aD} = (+l_3 + e\sigma)\,\hat{x}_b + h\,\hat{z}_b \\
    \vec{aA} = -l_2\,\hat{x}_b + (h - r_w)\,\hat{z}_b
\end{gather}

\subsubsection{Equations of motion}

% COMMENT: \dot{H}_a is the rate of change of angular momentum about the system CG a.
% It is not simply I_1*theta_ddot_1 because the spinning wheels and crank contribute;
% the expansion into individual inertia terms is shown in Rotation below.
The moment equation uses $\dot{\vec{H}}_a$, the rate of change of angular momentum about $a$, rather than a single body inertia, because the spinning wheels and crankset each contribute angular momentum independently:
\begin{align}
    m\,\vec{\ddot{x}}_1 &= \sum \vec{F}_1 = \vec{W}_1 + \vec{R}_{2,0C} + \vec{R}_{3,0D} + \vec{A}_{1f} \\[6pt]
    \dot{\vec{H}}_a &= \sum \vec{M}_1(a) = \vec{aC}\times\vec{R}_{2,0C} + \vec{aD}\times\vec{R}_{3,0D} + \vec{aA}\times\vec{A}_{1f} + \vec{A}_{1m}
\end{align}

\subsubsection{Translation}

% COMMENT: the vertical equation gives the static load constraint (ddot{z}_1 = 0, A_{1fz} = 0).
% It is used in Assembly to solve for normal reactions without needing the moment equation.
Weight is projected onto the road frame using the gradient angle $\theta_1$. The vertical equation enforces the constraint that the bike remains on the road:
\begin{align}
    m\,\ddot{x}_1 &= -mg\sin\theta_1 + R_{2,0Cx} + R_{3,0Dx} - A_{1fx} \\
    0               &= mg\cos\theta_1 + R_{2,0Cz} + R_{3,0Dz}
\end{align}

\subsubsection{Rotation}

% COMMENT: \dot{H}_y = sum I_k * theta_ddot_k because all bodies translate together so the
% transport terms in H vanish, and sum m_k * a_{a_k} = 0 since a is the system CG.
% The frame does not rotate (theta_ddot_1 = 0) so only wheels and crank appear on the left.
% The right side expands the cross products using the geometry defined above.
Since all subsystems translate together, the transport terms in $\dot{H}_a$ vanish and only the spin inertias remain. The frame does not rotate ($\ddot{\theta}_1 = 0$), so:
\begin{align}
    \dot{H}_y &= I_2\ddot{\theta}_2 + I_3\ddot{\theta}_3 + I_4\ddot{\theta}_4 \\
    \dot{H}_y &= h\left(R_{2,0Cx} + R_{3,0Dx}\right) + l_2'\,R_{2,0Cz} - l_3'\,R_{3,0Dz}
               - \left(h - r_w\right)A_{1fx} - A_{1my}
\end{align}

% =====================================================================
\newpage
\subsection{Wheels}
% =====================================================================

% COMMENT: each wheel is treated as an independent rigid body rotating about its fixed axle.
% The axle reaction R_{k,1} (frame on wheel) passes through the axle and produces no moment
% about it, so it does not appear in the wheel rotation equations.
Each wheel is an independent rigid body. Moments are taken about the axle to eliminate the axle reaction, which passes through it and produces no moment.

\textbf{Assumptions}
\begin{itemize}
    \item Each wheel is a rigid body with uniform mass distribution --- CG at the axle
    \item Axle fixed relative to the frame --- no suspension travel
    \item Wheel spin axis is $\hat{y}_b$ --- no camber or toe
    \item Contact patch is a point --- no contact patch width or length effects
    \item No lateral forces --- pure longitudinal contact resultant in the $x$--$z$ plane
    \item Brake force acts at a single effective radius $r_{bk,k}$ --- caliper geometry collapsed
    \item No wheel flex or tyre sidewall compliance
\end{itemize}

\subsubsection{Geometry}

% COMMENT: AC = BD because both wheels share the same radius and rolling resistance coefficient.
% A and B are the rear and front axle locations. F is the cog tooth contact point on the rear wheel.
% J and K are the brake pad contact points on rear and front respectively.
Both wheels share the same radius $r_w$ and rolling resistance coefficient $c_{rr}$, so the axle-to-contact vector is identical for both. The chain and brake moment arms are specific to each wheel:
\begin{equation}
    |\vec{AF}| = r_{cg} \qquad |\vec{AJ}| = r_{bk,2} \qquad |\vec{BK}| = r_{bk,3}
\end{equation}
The axle-to-contact resultant vector, expressed in the road frame:
\begin{equation}
    \vec{AC} = \vec{BD} = e\sigma\,\hat{x}_b + r_w\,\hat{z}_b
\end{equation}

\subsubsection{Forces}

% COMMENT: T_{2,4} pulls on the rear wheel at F along the chain span toward the chainring.
% B_{k,1} is the brake pad friction acting tangentially at the effective brake radius, opposing spin.
% The front wheel carries no chain force.
The chain tension acts only on the rear wheel at the cog contact point $F$. Brake forces act tangentially at the effective pad radius on whichever wheel is braked:
\begin{gather}
    \vec{T}_{2,4} = T_{2,4}\,\hat{e}_{span} \\
    \vec{B}_{k,1} = B_{k,1}\,\hat{e}_{t,bk}, \quad k = 2,3
\end{gather}

\subsubsection{Equations of motion}

% COMMENT: wheel weight passes through the axle (CG at axle for a uniform wheel), so it
% produces no moment about the axle either.
Expanding the cross products with the geometry above, noting that $\vec{R}_{k,1A}$ and $\vec{W}_k$ both act through the axle and vanish:
\begin{align}
    I_2\ddot{\theta}_2 &= \sum \vec{M}_2(A)\cdot\hat{y}
        = \bigl[\vec{AC}\times\vec{R}_{2,0C} + \vec{AF}\times\vec{T}_{2,4}
          + \vec{AJ}\times\vec{B}_{2,1}\bigr]\cdot\hat{y} \\[6pt]
    I_3\ddot{\theta}_3 &= \sum \vec{M}_3(B)\cdot\hat{y}
        = \bigl[\vec{BD}\times\vec{R}_{3,0D}
          + \vec{BK}\times\vec{B}_{3,1}\bigr]\cdot\hat{y}
\end{align}

\subsubsection{Rear wheel}

The traction force $R_{2,0Cx}$ drives the wheel; the normal load $R_{2,0Cz}$ contributes a rolling resistance moment through the offset $e\sigma$; chain tension resists spin; braking resists spin:
\begin{equation}
    I_2\ddot{\theta}_2 = r_w\,R_{2,0Cx} - e\sigma\,R_{2,0Cz}
        - r_{cg}\,T_{2,4} + \sigma\,r_{bk,2}\,B_{2,1}
\end{equation}

\subsubsection{Front wheel}

The front wheel carries no chain force. Only the contact resultant and the brake act:
\begin{equation}
    I_3\ddot{\theta}_3 = r_w\,R_{3,0Dx} - e\sigma\,R_{3,0Dz}
        + \sigma\,r_{bk,3}\,B_{3,1}
\end{equation}

% =====================================================================
\newpage
\subsection{Drivetrain}
% =====================================================================

% COMMENT: the drivetrain here is the crankset only (chainring + cranks + pedals).
% The cassette is part of the rear wheel subsystem. A fixed gear ratio r_cg/r_cr is assumed.
The drivetrain subsystem comprises the crankset: chainring, cranks, and pedals. The cassette is treated as part of the rear wheel. A fixed gear ratio $r_{cg}/r_{cr}$ is assumed throughout.

\textbf{Assumptions}
\begin{itemize}
    \item Fixed gear ratio --- no derailleur shifting, single chainring and cog
    \item No drivetrain efficiency loss --- chain transmits torque without friction
    \item Chain mass neglected --- chain inertia not modelled
    \item Crankset is a rigid body with CG at the BB axle $E$
    \item Pedal force acts tangentially at the crank tip --- no radial component
    \item Left and right cranks combined as a single effective torque
    \item No freewheel modelled explicitly --- $T_{2,4} \ge 0$ enforced as a post-check
\end{itemize}

\subsubsection{Geometry}

% COMMENT: E is the bottom bracket axle. G is the chainring tooth contact point.
% H_L, H_R are the left and right pedal spindle locations; r_p is the crank length.
% Moments are taken about E to eliminate the BB axle reaction and crankset weight.
$E$ is the bottom bracket axle, the rotation centre of the crankset. The chainring engages the chain at $G$; the rider applies force at the pedal spindles $H_L$, $H_R$. The relevant moment arms are:
\begin{equation}
    |\vec{EG}| = r_{cr} \qquad |\vec{EH_L}| = |\vec{EH_R}| = r_p
\end{equation}

\subsubsection{Forces}

% COMMENT: T_{4,2} is the Newton's third law reaction to T_{2,4}: the chain pulls the chainring
% forward (driving direction) with the same magnitude it pulls the cog backward.
% P_{4,1s} is the pedal force from the rider, signed positive in the driving direction.
The chain pulls the chainring at $G$ with the same tension it pulls the cog, by Newton's third law. The rider applies pedal forces tangentially at each crank tip:
\begin{gather}
    \vec{T}_{4,2} = -\vec{T}_{2,4} \\
    \vec{P}_{4,1s} = P_{4,1s}\,\hat{e}_{t,s}, \quad s = L,R
\end{gather}

\subsubsection{Equations of motion}

% COMMENT: BB axle reaction and crankset weight both pass through E and produce no moment.
Moments about $E$ eliminate the BB axle reaction and crankset weight:
\begin{equation}
    I_4\ddot{\theta}_4 = \sum \vec{M}_4(E)\cdot\hat{y}
        = \bigl[\vec{EG}\times\vec{T}_{4,2}
          + \vec{EH_L}\times\vec{P}_{4,1L}
          + \vec{EH_R}\times\vec{P}_{4,1R}\bigr]\cdot\hat{y}
\end{equation}

\subsubsection{Crankset}

Pedal force drives the crankset; chain tension resists it:
\begin{equation}
    I_4\ddot{\theta}_4 = r_{cr}\,T_{2,4} - r_p\left(P_{4,1L} + P_{4,1R}\right)
\end{equation}

% =====================================================================
\newpage
\subsection{Assembly}
% =====================================================================

% COMMENT: the four subsystems are coupled through three kinematic constraints: no-slip at both
% contacts and the fixed chain ratio. These reduce all angular accelerations to one unknown: ddot{x}_1.
The four subsystems are coupled through three kinematic constraints: no-slip at each contact and the fixed chain ratio. All angular accelerations reduce to a single unknown $\ddot{x}_1$. The solve order has no algebraic loop: the EOM gives $\ddot{x}_1$ first, then reactions are post-processed in sequence.

\textbf{Assumptions}
\begin{itemize}
    \item No-slip at both contacts --- wheels roll without sliding
    \item Kinematic constraint is holonomic --- $\dot{x}_1 = -r_w\dot{\theta}_2 = -r_w\dot{\theta}_3$ at all times
    \item Chain does not slip on sprockets --- gear ratio is exact
    \item All assumptions from Bike, Wheels, and Drivetrain carry forward
\end{itemize}

\subsubsection{Kinematics}

% COMMENT: no-slip gives theta_ddot_k = -x_ddot / r_w. The sign follows from hat{y} pointing
% right: forward motion (positive x_dot) corresponds to negative wheel rotation.
% The chain ratio r_cg/r_cr then fixes theta_ddot_4.
No-slip at each contact relates wheel angular acceleration to frame linear acceleration. The chain ratio propagates this to the crankset:
\begin{align}
    \ddot{\theta}_1 &= 0 \\
    \ddot{\theta}_2 = \ddot{\theta}_3 &= -\frac{\ddot{x}_1}{r_w} \\
    \ddot{\theta}_4 &= -\frac{r_{cg}}{r_{cr}\,r_w}\,\ddot{x}_1
\end{align}

\subsubsection{Chain tension}

% COMMENT: kinematics substituted into the crankset equation and solved for T_{2,4}.
% The first term is inertial (crankset resists acceleration); the second is the pedal input.
% If T_{2,4} < 0 the chain would need to push, which it cannot: freewheel condition.
Substituting kinematics into the crankset equation and solving for $T_{2,4}$. The first term captures crankset inertia; the second is the pedal input referred to the chain:
\begin{equation}
    \boxed{T_{2,4} = -\frac{I_4\,r_{cg}}{r_{cr}^2\,r_w}\,\ddot{x}_1
           + \frac{r_p}{r_{cr}}\left(P_{4,1L} + P_{4,1R}\right)}
\end{equation}

\subsubsection{Traction}

% COMMENT: kinematics substituted into each wheel equation and solved for R_{k,0x}.
% These are the horizontal ground forces required to sustain the motion — checked against
% the traction limit in Limits.
Substituting kinematics into each wheel equation and solving for the longitudinal ground reaction. These are the horizontal forces the ground must supply; they are checked against the friction limit below:
\begin{align}
    \boxed{R_{2,0Cx} = -\frac{I_2}{r_w^2}\,\ddot{x}_1
        + c_{rr}\sigma\,R_{2,0Cz}
        + \frac{r_{cg}}{r_w}\,T_{2,4}
        - \frac{\sigma\,r_{bk,2}}{r_w}\,B_{2,1}} \\
    \boxed{R_{3,0Dx} = -\frac{I_3}{r_w^2}\,\ddot{x}_1
        + c_{rr}\sigma\,R_{3,0Dz}
        - \frac{\sigma\,r_{bk,3}}{r_w}\,B_{3,1}}
\end{align}

\subsubsection{Normal reactions}

% COMMENT: obtained from the bike rotation equation after substituting kinematics and traction.
% K collects the moment terms that do not depend on the front/rear split.
% The EOM itself needs only the sum R_{2,0Cz} + R_{3,0Dz} = -mg cos(theta), so the split
% is pure post-processing and does not affect the equation of motion.
The rotation equation, after substituting kinematics and traction, yields the auxiliary quantity $K$ which collects all contributions to the pitching moment. The normal reactions follow from $K$ and the vertical equilibrium. Note that the EOM needs only the sum of the normal loads, so this step is post-processing:
\begin{align}
    \boxed{R_{2,0Cz} = \frac{K - l_3'\,mg\cos\theta_1}{L}} \qquad
    \boxed{R_{3,0Dz} = -\frac{K + l_2'\,mg\cos\theta_1}{L}}
\end{align}
\begin{equation}
    K = -\frac{J}{r_w}\,\ddot{x}_1
        - hm\!\left(\ddot{x}_1 + g\sin\theta_1\right)
        - r_w\,A_{1fx} + A_{1my},
    \qquad
    J = I_2 + I_3 + \frac{r_{cg}}{r_{cr}}\,I_4
\end{equation}

\subsubsection{Equation of motion}

% COMMENT: obtained by substituting the traction expressions into the bike translation equation
% and using R_{2,0Cz} + R_{3,0Dz} = -mg cos(theta) to collapse the rolling resistance terms.
% m_eq is the total inertia of the system referred to the linear acceleration of the frame;
% it includes the rotational inertia of both wheels and the crankset reflected through the gear ratio.
Substituting traction into the bike translation equation and using $R_{2,0Cz} + R_{3,0Dz} = -mg\cos\theta_1$ collapses the rolling resistance terms. The effective mass $m_{eq}$ reflects all rotational inertias to the frame linear acceleration:
\begin{equation}
    \resizebox{\linewidth}{!}{$\displaystyle
    \boxed{
    m_{eq}\,\ddot{x}_1 =
        \frac{r_p\,r_{cg}}{r_{cr}\,r_w}\left(P_{4,1L} + P_{4,1R}\right)
        - \frac{\sigma}{r_w}\left(r_{bk,2}\,B_{2,1} + r_{bk,3}\,B_{3,1}\right)
        - mg\sin\theta_1
        - \sigma\,c_{rr}\,mg\cos\theta_1
        - A_{1fx}
    }$}
\end{equation}
\begin{equation}
    m_{eq} = m + \frac{I_2 + I_3}{r_w^2} + \frac{I_4\,r_{cg}^2}{r_{cr}^2\,r_w^2}
\end{equation}

\subsubsection{Limits}

The following conditions must hold for the model assumptions to remain valid. Violation of any one signals a regime change: wheel lift, chain slack, or sliding contact:
\begin{align}
    R_{2,0Cz} &\le 0, \qquad R_{3,0Dz} \le 0 \qquad \text{(no wheel lift)} \\
    \left|R_{2,0Cx}\right| &\le -\mu_s\,R_{2,0Cz}, \qquad
    \left|R_{3,0Dx}\right| \le -\mu_s\,R_{3,0Dz} \qquad \text{(no slip)} \\
    T_{2,4} &\ge 0 \qquad \text{(no freewheel)}
\end{align}

% =====================================================================
\newpage
\subsection{Analysis}
% =====================================================================

\textbf{Assumptions}
\begin{itemize}
    \item Steady state --- $\ddot{x}_1 = 0$, constant speed on constant gradient
    \item No braking --- $B_{2,1} = B_{3,1} = 0$
    \item Forward motion --- $\sigma = 1$
    \item Still air or head/tailwind only --- drag enters as $\tfrac{1}{2}\rho c_x A v_{rel}^2$
    \item Drivetrain efficiency not modelled --- power input equals power at the chain
    \item All limits satisfied --- no slip, no wheel lift, no freewheel
\end{itemize}

\subsubsection{Equilibrium}

% COMMENT: steady state on a constant gradient, no braking, forward motion (sigma = 1).
% Expressing the left side as power P eliminates the gear ratio and crank geometry —
% only drivetrain efficiency (not yet modelled) would remain if added.
Setting $\ddot{x}_1 = 0$, $B_{2,1} = B_{3,1} = 0$, $\sigma = 1$ gives the steady-state force balance. Three resistance terms appear: gradient, rolling, and aerodynamic drag, scaling respectively as $v^0$, $v^0$, and $v^2$:
\begin{equation}
    \frac{r_p\,r_{cg}}{r_{cr}\,r_w}\left(P_{4,1L} + P_{4,1R}\right) = mg\left(c_{rr}\cos\theta_1 + \sin\theta_1\right) + \tfrac{1}{2}\rho\,c_x A\,v^2
\end{equation}
Expressing the left side as prescribed power $\mathcal{P}$ using $|\dot{\theta}_4| = r_{cg}v/(r_{cr}r_w)$ eliminates the gear geometry:
\begin{equation}
    \frac{\mathcal{P}}{v} = mg\left(c_{rr}\cos\theta_1 + \sin\theta_1\right) + \tfrac{1}{2}\rho\,c_x A\,v^2
\end{equation}

\subsubsection{Sensitivity}

% COMMENT: full treatment following the uploaded FX27 document and Notion §6 analysis targets.
% The equilibrium cubic is simplified to P = mgv sin(theta) + kv^3 by dropping rolling
% resistance (order of magnitude smaller) to keep the sensitivity structure clean.
% Elasticities are derived by implicit differentiation; three questions are then answered
% by reading off the two gradient limits.

Write the aerodynamic drag power as $k \equiv \tfrac{1}{2}\rho\,c_x A$ and drop rolling resistance
(an order of magnitude below gravity and drag at any sustained speed), leaving the simplest model
that still contains both effects each question is about:
\begin{equation}
    \mathcal{P} = mg\,v\sin\theta_1 + k\,v^3
\end{equation}

\paragraph{Elasticities.}
The elasticity $E_p \equiv (p/v)(\partial v/\partial p)$ gives the fractional change in speed per
fractional change in parameter $p$, at constant $\mathcal{P}$. Write the balance as
$G(v;\,m,k,\theta_1,\mathcal{P}) = mg\,v\sin\theta_1 + k\,v^3 - \mathcal{P} = 0$ and
differentiate implicitly:
\begin{equation}
    \frac{\partial v}{\partial p} = -\frac{\partial G/\partial p}{\partial G/\partial v},
    \qquad
    \frac{\partial G}{\partial v} = mg\sin\theta_1 + 3kv^2 \equiv D
\end{equation}
\begin{equation}
    \boxed{E_m = -\frac{mg\sin\theta_1}{D}, \qquad
           E_k = -\frac{kv^2}{D}, \qquad
           E_{\mathcal{P}} = \frac{mg\sin\theta_1 + kv^2}{D}}
\end{equation}
Since $G$ is linear in $(m, k, \mathcal{P})$ at fixed $v$, Euler's theorem gives the exact identity:
\begin{equation}
    E_m + E_k + E_{\mathcal{P}} = 0
\end{equation}

\paragraph{Two gradient limits.}
\textbf{Flat, aero-dominated} ($mg\sin\theta_1 \ll kv^2$, $D \approx 3kv^2$):
\begin{equation}
    E_m \approx 0, \qquad E_k \approx -\tfrac{1}{3}, \qquad E_{\mathcal{P}} \approx \tfrac{1}{3}
    \quad\Longrightarrow\quad v \propto \mathcal{P}^{1/3},\;\; v \propto k^{-1/3}
\end{equation}
Weight is irrelevant on the flat; speed depends only on power and drag, and weakly (cube root).

\textbf{Steep, gravity-dominated} ($kv^2 \ll mg\sin\theta_1$, $D \approx mg\sin\theta_1$):
\begin{equation}
    E_m \approx -1, \qquad E_k \approx 0, \qquad E_{\mathcal{P}} \approx 1
    \quad\Longrightarrow\quad v \approx \frac{\mathcal{P}}{mg\sin\theta_1}
\end{equation}
Aerodynamics is irrelevant while climbing; speed is set entirely by power-to-weight.

\paragraph{Does aerodynamics matter more for a lighter rider?}
Two riders with the same $\mathcal{P}$ and $k$, differing only in mass. Since $E_m < 0$ the
lighter rider is faster. Differentiate $|E_k| = kv^2/D$ along the equilibrium curve:
\begin{equation}
    \frac{\mathrm{d}|E_k|}{\mathrm{d}v} = \frac{2kv\,mg\sin\theta_1}{D^2} > 0
\end{equation}
$|E_k|$ is strictly increasing in $v$, so aerodynamics matters strictly more to the lighter
(faster) rider. The gap is smallest while climbing ($|E_k|\approx 0$ for both) and largest on
the flat ($|E_k|\to\tfrac{1}{3}$ for both).

\paragraph{Is there a best rider weight?}
Comparing riders of different size: power and frontal area both scale with mass. Model with
exponents against a reference rider $(m_0, \mathcal{P}_0, k_0)$:
\begin{equation}
    \mathcal{P}(m) = \mathcal{P}_0\!\left(\frac{m}{m_0}\right)^{a}, \qquad
    k(m) = k_0\!\left(\frac{m}{m_0}\right)^{b}
\end{equation}
($b \approx \tfrac{2}{3}$ from geometric similarity; $a$ is physiological, typically $a < 1$.)
The size elasticity follows by chain rule:
\begin{equation}
    \boxed{\left.\frac{\mathrm{d}\ln v}{\mathrm{d}\ln m}\right|_{\text{across riders}}
    = E_m + a\,E_{\mathcal{P}} + b\,E_k}
\end{equation}
Plugging in the two limits: climbing gives $\approx a - 1 < 0$ (smaller always wins for $a < 1$);
flat gives $\approx (a-b)/3$ (bigger wins if power scales faster than frontal area, $a > b$).
Best weight flips sign with gradient; the crossover depends on $a - b$, not on mass directly.

\paragraph{Is a steeper climb more energy-efficient for the same height gain?}
Fix height $\Delta h$; route length is $d(\gamma) = \Delta h / \sin\theta_1$. Energy spent at
constant power $\mathcal{P}$:
\begin{equation}
    E(\theta_1) = \mathcal{P}\,\frac{d(\theta_1)}{v(\theta_1)}
    = \frac{\mathcal{P}\,\Delta h}{v(\theta_1)\sin\theta_1}
\end{equation}
From the power balance, $v\sin\theta_1 = (\mathcal{P} - kv^3)/(mg)$, which is maximised by
minimising $v$. Since $\partial v/\partial\theta_1 < 0$ (steeper is slower),
$v\sin\theta_1$ rises monotonically with gradient and $E(\theta_1)$ falls:
\begin{equation}
    \frac{\mathrm{d}E}{\mathrm{d}\theta_1} < 0
    \;\Longrightarrow\;
    \text{steeper is always more energy-efficient for a fixed height gain}
\end{equation}
Mechanism: a steeper route is shorter and slower; aerodynamic loss scales as $v^2 \times
\text{distance}$, so cutting speed helps cubically while the shorter distance helps linearly.

\subsubsection{Stoppie limit}

% COMMENT: rear wheel lift under hard front braking. Setting R_{2,0Cz} = 0 in the normal
% reaction expression gives K directly; substituting into the K definition yields the
% critical deceleration. Aero terms are retained for completeness.
Front braking transfers load forward. The rear wheel lifts when $R_{2,0Cz} = 0$. Setting this in the normal reaction expression fixes $K$, which is then substituted into its definition to give the critical deceleration:
\begin{equation}
    R_{2,0Cz} = 0 \;\Rightarrow\; K = l_3'\,mg\cos\theta_1
\end{equation}
\begin{equation}
    \ddot{x}_{1,\text{stoppie}} = \frac{-\left(l_3'\,mg\cos\theta_1 + hm\,g\sin\theta_1 + r_w\,A_{1fx} - A_{1my}\right)}{J/r_w + hm}
\end{equation}

\subsubsection{Wheelie limit}

% COMMENT: front wheel lift under hard acceleration. Same procedure with R_{3,0Dz} = 0.
% On a steep climb the gradient term assists the wheelie; the aero pitching moment A_{1my}
% acts nose-down and therefore resists it.
Hard acceleration transfers load rearward. The front wheel lifts when $R_{3,0Dz} = 0$:
\begin{equation}
    R_{3,0Dz} = 0 \;\Rightarrow\; K = -l_2'\,mg\cos\theta_1
\end{equation}
\begin{equation}
    \ddot{x}_{1,\text{wheelie}} = \frac{l_2'\,mg\cos\theta_1 - hm\,g\sin\theta_1 - r_w\,A_{1fx} + A_{1my}}{J/r_w + hm}
\end{equation}

\subsubsection{Traction limit}

% COMMENT: slip is checked per wheel after the normal reactions are known from post-processing.
% Under the no-slip assumption it does not feed back into the EOM — it is a validity check only.
% The available friction budget is split between driving and braking at each wheel; exceeding
% it at either wheel invalidates the model for that wheel.
Slip is checked after the normal reactions are known. It does not feed back into the EOM — the no-slip assumption is either valid or it is not, and violation signals that a different model (e.g.\ with combined slip) is required. The available friction at each contact is $\mu_s |R_{k,0z}|$; the longitudinal force must stay within it:
\begin{align}
    \left|R_{2,0Cx}\right| &\le -\mu_s\,R_{2,0Cz} \;\Rightarrow\; \text{rear contact valid} \\
    \left|R_{3,0Dx}\right| &\le -\mu_s\,R_{3,0Dz} \;\Rightarrow\; \text{front contact valid}
\end{align}

\end{document}